\bookchapter{ملحق الفصل 23: دليل البدء في GitHub — من استضافة الأكواد إلى منصة التعاون}

\chapterintro{في الفصل السابق تحدثنا عن “التفكير البرمجي”، ولترجمة هذا التفكير إلى واقع والعمل مع آخرين في مشاريع، لا مفر من اسم — GitHub. سواء كنت طالب علوم حاسوب تشارك في مسابقات، أو تريد حفظ أعمالك البرمجية، أو تريد ببساطة رؤية ما يفعله المطورون حول العالم، فإن GitHub منصة يجب أن تعرفها. في هذا القسم سنبدأ من المفاهيم الأساسية لمساعدتك في بناء فهم واضح.}

\booksection{أولًا: ما هو GitHub بالضبط؟}

Git\textbf{GitHub}الجوهر

أداة تحكم في الإصدارات (برنامج)\textbf{منصة استضافة الأكواد (موقع ويب)}موقع التشغيل

على الحاسوب المحلي

\booksection{ثانيًا: الفرق بين Git و GitHub (مفاهيم أساسية!)}

المسؤولية الأساسية

\begin{center}
\begin{tabularx}{\linewidth}{|X|X|X|}
\hline
\textbf{تسجيل كل تغيير في الملفات (كيف تغير، من غيّر، متى)} & \textbf{وضع مستودعات Git في السحابة لتسهيل التخزين والعرض والتعاون} & \textbf{هل يحتاج إلى إنترنت} \tabularnewline
\hline
لا يحتاج، يعمل دون اتصال & يحتاج إلى اتصال للوصول & الملكية \tabularnewline
برنامج مفتوح المصدر، أنشأه Linus Torvalds (أبو Linux) & منتج مملوك لـ Microsoft & منافسون من نفس الفئة \tabularnewline
SVN وMercurial & GitLab وGitee (مايون) وBitbucket & ببساطة: \tabularnewline
Git يدير "كيفية تسجيل التغييرات"، و GitHub يدير "أين توضع التغييرات وكيف نتعاون" & . Git أداة تعمل محليًا؛ و GitHub منصة تضع مستودعات Git في السحابة لتسهيل التعاون والعرض. يمكنك استخدام Git فقط دون GitHub (مثل كتابة الأكواد وإدارة إصداراتها بنفسك)، لكن استخدام GitHub يتطلب استخدام Git (لأن GitHub مبني على Git). & ثالثًا: المفاهيم والوظائف الأساسية في GitHub \tabularnewline
مفاهيم أساسية & قبل البدء، تعرف على بعض الكلمات المتكررة، وإلا ستكون واجهة GitHub محيرة: & المستودع (Repository، اختصارًا Repo) \tabularnewline
: كل مشروع يقابله مستودع، يحتوي على جميع ملفات الأكواد لهذا المشروع والإصدارات التاريخية والوثائق. يمكن فهمه على أنه "مجلد المشروع". & الالتزام (Commit) & : سجل تعديل واحد على الكود. يرافق كل Commit رسالة توضيحية تسجل "ما الذي تغير هذه المرة، ولماذا". هذا هو أساس قدرة Git/GitHub على "السفر عبر الزمن". \tabularnewline
\hline
\end{tabularx}
\end{center}

الفرع (Branch)\textbf{: خط تطوير مستقل منفصل عن المسار الرئيسي. على سبيل المثال، إذا أردت إضافة ميزة جديدة، يمكنك فتح فرع للتطوير المنفصل، ثم الدمج مع المسار الرئيسي بعد الانتهاء، دون التأثير على الآخرين.}الاستنساخ (Clone)

\booksection{ثالثًا: المفاهيم والوظائف الأساسية في GitHub}

\booksubsection{مفاهيم أساسية}

: Push هو رفع الكود المحلي إلى GitHub السحابية، و Pull هو تنزيل أحدث كود من GitHub إلى المحلي. هذان هما الإجراءان الأساسيان لمزامنة الكود أثناء عمل الفريق.

\begin{itemize}
  \item \textbf{طلب الدمج (Pull Request، اختصارًا PR)}: عندما تنتهي من تعديل الكود على فرع وتريد دمجه في الفرع الرئيسي، تطلق PR ليراجع أعضاء الفريق الكود قبل الدمج. هذه هي آلية التعاون الأساسية في GitHub.
  \item \textbf{المسألة (Issue)}: تُستخدم لتسجيل الأخطاء وطلبات الميزات ومواضيع النقاش. كل Issue مثل منشور نقاش، يمكن تعيين مسؤول له ووضع علامات.
  \item \textbf{Fork}: نسخ مستودع شخص آخر إلى حسابك، يمكنك التعديل بحرية على نسختك دون التأثير على المستودع الأصلي. بعد الانتهاء يمكنك إرسال PR للمؤلف الأصلي لطلب دمج تحسيناتك.
  \item \textbf{الوظائف الأساسية}ما يقدمه GitHub ليس مجرد تخزين للأكواد، بل يشمل أيضًا:
  \item \textbf{تحكم الإصدارات}: تسجيل تاريخ كل تغيير في الكود، والعودة في أي وقت إلى آخر إصدار سليم عند تلف الكود.
  \item \textbf{تعاون الفريق}: دعم التطوير المتوازي لنفس المشروع من عدة مستخدمين، دون الكتابة فوق بعضهم البعض أو التعارض.
  \item \textbf{تتبع المشكلات}: استخدام Issue لتسجيل الأخطاء وطلبات الميزات، لتسهيل متابعة الفريق.
  \item \textbf{إدارة المشاريع}: تكامل أدوات إدارة مثل المهام وألواح المشروع (Projects)، أي ما يعادل Trello خفيف مدمج.
\end{itemize}

\booksubsection{الوظائف الأساسية}

: عبر

\begin{itemize}
  \item \textbf{GitHub Actions}لتنفيذ التجميع والاختبار والتعبئة والنشر تلقائيًا بعد إرسال الكود، وهي أداة أساسية في "CI/CD (التكامل المستمر/النشر المستمر)" الحديث.
  \item \textbf{إدارة الوثائق}: README يقدم المشروع، و Wiki يخزن البنية المعمارية التفصيلية ومعايير التطوير، ويمكن لـ Pages استضافة مواقع ثابتة مباشرة.
  \item \textbf{استضافة الأكواد وعرضها}: لكل مستودع صفحة ويب جميلة العرض، مع دعم تلوين الصيغة، وهي "واجهة" أعمالك التقنية.
  \item \textbf{مجتمع مفتوح المصدر}: يمكنك Star (مفضلة) المشاريع التي تعجبك، و Watch (متابعة) مستجدات المشاريع، و Follow (متابعة) المطورين، والمشاركة في التعاون مفتوح المصدر عالميًا.
  \item \textbf{رابعًا: التسجيل والبدء — أنشئ أول مستودع لك في خمس دقائق}تسجيل حساب\textbf{زر}github.com
  \item \textbf{، وانقر على "Sign up".}املأ البريد الإلكتروني، واضبط كلمة المرور واسم المستخدم (سيظهر اسم المستخدم في عنوان صفحتك الشخصية، يُنصح بالإنجليزية، مثل
  \item \textbf{github.com/yourname}).
  \item \textbf{أكمل التحقق من البريد الإلكتروني والتحقق البشري.}اختر الخطة المجانية (Free)، وهي كافية تمامًا للمستخدمين الأفراد.
\end{itemize}

\booksection{رابعًا: التسجيل والبدء — أنشئ أول مستودع لك في خمس دقائق}

\booksubsection{تسجيل حساب}

\begin{enumerate}
  \item املأ اسم المستودع (Repository name)، مثل\href{https://github.com}{my-first-project}.
  \item اكتب وصفًا موجزًا (Description) ليعرف الآخرون ما يفعله هذا المستودع.\texttt{اختر عام (Public) أو خاص (Private).\allowbreak  المستودعات العامة يراها الجميع، والخاصة يراها أنت والمتعاونون فقط.\allowbreak  الحسابات المجانية تدعم أيضًا مستودعات خاصة غير محدودة.\allowbreak }حدد "Add a README file" — README هو "كتيب التعليمات" للمشروع، ويُنصح بوجوده في كل مستودع.
  \item انقر على "Create repository"، وقد وُلد أول مستودع GitHub لك!
  \item معلومة صغيرة: ما هو README
\end{enumerate}

\booksubsection{إنشاء أول مستودع}

\begin{enumerate}
  \item خامسًا: سير العمل الأساسي — من المحلي إلى السحابة
  \item لفهم سير عمل GitHub، المفتاح هو تذكر هذه الدورة:\texttt{تعديل محلي تحضير التزام دفع تحديث سحابي}.
  \item تعديل الملفات
  \item : كتابة الكود أو تعديل المستندات على الحاسوب المحلي باستخدام محرر.
  \item git add
  \item : إضافة الملفات المعدلة إلى "منطقة التحضير"، أي إخبار Git "هذه التغييرات أستعد لتسجيلها".
\end{enumerate}

\begin{tcbinfo}{git commit}
: تسجيل تغييرات منطقة التحضير رسميًا كـ “التزام”، مع رسالة توضيحية، مثل “إصلاح خطأ تسجيل الدخول”. تُنجز هذه الخطوة محليًا فقط.

\end{tcbinfo}

\booksection{خامسًا: سير العمل الأساسي — من المحلي إلى السحابة}

: رفع الالتزامات المحلية إلى GitHub السحابية، ليتمكن الآخرون من رؤيتها.\textbf{git pull}: قبل بدء العمل، اسحب أحدث كود من السحابة لضمان مزامنة المحلي مع بقية الفريق.

\begin{enumerate}
  \item \textbf{للمبتدئين، تذكر هذه القاعدة:}"اسحب أولًا ثم عدّل، بعد التعديل add، و commit مع رسالة، وأخيرًا push"
  \item \textbf{. اكتسب هذه العادة ولن يختلط عمل الفريق.}سادسًا: سير عمل تعاون الفريق — كيف يعمل Pull Request
  \item \textbf{أقوى جانب في GitHub هو آلية التعاون. سير عمل التعاون الفريقي النموذجي هو:}فتح فرع
  \item \textbf{: افتح فرعًا جديدًا من الفرع الرئيسي (يُسمى عادة main أو master)، مثل}feature/login
  \item \textbf{، وطوّر الميزة الجديدة على هذا الفرع.}التطوير والالتزام
\end{enumerate}

: اكتب الكود على الفرع الجديد، و commit عدة مرات لتسجيل التقدم.\textbf{الدفع إلى GitHub}: ارفع الفرع إلى المستودع السحابي.

\booksection{سادسًا: سير عمل تعاون الفريق — كيف يعمل Pull Request}

: أطلق PR على صفحة GitHub، لطلب دمج فرعك في الفرع الرئيسي. اكتب بوضوح ما الذي تغير في هذا PR ولماذا.

\begin{enumerate}
  \item \textbf{مراجعة الكود (Code Review)}: يطّلع أعضاء الفريق على كودك ويطرحون آراء أو اقتراحات للتعديل. يمكنك الاستمرار في commit بناءً على التغذية الراجعة، وسيُحدَّث PR تلقائيًا.\texttt{الدمج}: بعد اجتياز المراجعة، انقر "Merge" لدمج الكود في الفرع الرئيسي. وقد دخلت مساهمتك المشروع رسميًا!
  \item \textbf{يضمن هذا التدفق أن}أي كود يدخل الفرع الرئيسي يمر بالمراجعة
  \item \textbf{، مما يقلل بشكل كبير من احتمال الخطأ. المشاريع مفتوحة المصدر تستخدم أيضًا هذه الآلية لقبول مساهمات الغرباء حول العالم — يمكنك تمامًا تقديم PR في مشروع مفتوح المصدر معروف، وبعد قبوله تصبح مساهمًا في هذا المشروع.}سابعًا: سيناريوهات التطبيق النموذجية لـ GitHub
  \item \textbf{مسابقات الطلاب ومشاريع المقررات}يستخدم طلاب علوم الحاسوب GitHub في مسابقات ACM وهاكاثون ومشاريع التخرج. فهو منصة تعاون فريقي ونافذة عرض للأعمال — يمكن لجنة التحكيم أن تطّلع مباشرة على مستودعك لتفهم جودة الكود وعملية التعاون. كثير من الشركات تطلب أيضًا إرفاق صفحة GitHub الرئيسية عند التوظيف، كدليل على القدرة التقنية.
  \item \textbf{قاعدة المعرفة الشخصية واستضافة الملاحظات}غير المبرمجين يستطيعون استخدام GitHub أيضًا. كثيرون يستخدمونه لاستضافة ملاحظات التعلم والمدونات التقنية (مع GitHub Pages يمكن إنشاء موقع مجاني) والكتب الإلكترونية. الكتابة بتنسيق Markdown + إدارة الإصدارات بـ Git + النسخ الاحتياطي السحابي، ثلاثة في واحد.
  \item \textbf{المشاركة في المشاريع مفتوحة المصدر}GitHub هو المعقل الأساسي لحركة مفتوح المصدر عالميًا. يمكنك Fork المشاريع التي تعجبك، وتعديلها ثم إرسال PR لرد الجميل للمجتمع؛ أو ببساطة Star للمشاريع الجيدة ومتابعة التحديثات. المشاركة في مفتوح المصدر طريق ممتاز لتحسين التقنية وبناء السمعة.
\end{enumerate}

النشر الآلي و CI/CD\textbf{بمساعدة GitHub Actions، يمكن تحقيق تدفق آلي بالكامل “دفع الكود، اختبار تلقائي + نشر تلقائي على الخادم”. هكذا يُبنى كثير من المدونات الشخصية ومواقع الوثائق، تكتب المقال وتدفعها فيتحدث الموقع تلقائيًا.}النسخ الاحتياطي للكود والمزامنة بين الأجهزة

\booksection{سابعًا: سيناريوهات التطبيق النموذجية لـ GitHub}

\booksubsection{مسابقات الطلاب ومشاريع المقررات}

اكتب رسائل الالتزام بوضوح

\booksubsection{قاعدة المعرفة الشخصية واستضافة الملاحظات}

التزم كثيرًا بخطوات صغيرة

\booksubsection{المشاركة في المشاريع مفتوحة المصدر}

طوّر على فرع، لا تعدّل الفرع الرئيسي مباشرة

\booksubsection{النشر الآلي و CI/CD}

أحسن استخدام .gitignore

\booksubsection{النسخ الاحتياطي للكود والمزامنة بين الأجهزة}

اطّلع على مشاريع مفتوحة المصدر الممتازة

\booksection{ثامنًا: نصائح للمبتدئين}

\begin{itemize}
  \item \textbf{تقدم بحزمة طالب GitHub}: بعد مصادقة الطلاب يمكنك الحصول مجانًا على حصص كبيرة من الأدوات والخدمات المدفوعة (Copilot، أسماء النطاقات، الخدمات السحابية، إلخ)، بقيمة آلاف الدولارات.
  \item \textbf{تذكير أمني}لا ترفع أبدًا معلومات حساسة مثل كلمات المرور ومفاتيح API وسلاسل اتصال قواعد البيانات إلى GitHub، خاصة في المستودعات العامة — بمجرد التسريب ستلتقطها برامج الزحف فورًا والعواقب وخيمة.
  \item \textbf{إذا رفعت معلومات حساسة بالخطأ، فحذف الملف لا يكفي (لأنها لا تزال في السجل التاريخي)، استخدم}git filter-branch
  \item \textbf{أو أدوات متخصصة لتنظيف التاريخ، وغيّر فورًا المفاتيح المسربة.}المستودعات الخاصة ليست آمنة تمامًا أيضًا، يُفضل إدارة المعلومات الحساسة بمتغيرات البيئة، ولا تدخلها في مستودع الكود.
  \item \textbf{ملاحظات للوصول إلى GitHub من داخل الصين}قد يكون الوصول إلى GitHub من داخل الصين غير مستقر أحيانًا، فقد لا يفتح أو يكون الاستنساخ بطيئًا. يمكنك تجربة: (1) استخدام مواقع مرآة محلية (مثل GitClone وghproxy) لتسريع الاستنساخ؛ (2) تكوين Hosts أو استخدام وكيل؛ (3) إذا كان التعاون موجّهًا أساسًا للداخل، يمكنك أيضًا التفكير في البديل الصيني
  \item \textbf{Gitee (مايون)}، بوظائف مماثلة وسرعة وصول عالية. كثير من الفرق تستضيف على GitHub و Gitee في الوقت ذاته، لجمع الفوائد الدولية والسرعة المحلية.
\end{itemize}

\begin{tcbwarn}{مستوى متقدم: استخدم GitHub Pages لبناء مدونة شخصية مجانية}
\begin{itemize}
  \item يسمح لك GitHub Pages باستضافة موقع ثابت مجانًا، وبالتعاون مع مولدات المواقع الثابتة مثل Jekyll وHexo وHugo، يمكنك بناء مدونة شخصية بسهولة. اكتب مقال Markdown وادفعه إلى المستودع، فتتحدث المدونة تلقائيًا، وتأتي مع HTTPS واسم نطاق مخصص. هكذا بُنيت مدونات عدد لا يحصى من المدونين التقنيين، بتكلفة صفرية وبدون إعلانات وإصدارات قابلة للتحكم.
  \item تاسعًا: الخلاصة — لماذا يجب على الجميع فهم GitHub\texttt{لم يعد GitHub مجرد "أداة المبرمجين"، بل أصبح أشبه بـ}طريقة عمل في العصر الرقمي
  \item : إدارة الإصدارات تجعل كل شيء قابلًا للتتبع، وآليات التعاون تجعل الفريق يعمل بكفاءة، وروح مفتوح المصدر تجعل المعرفة تتدفق بحرية. حتى لو لم تكتب كودًا، فإن فهم مفاهيم "الالتزام والفرع والدمج" يساعدك على إدارة أي محتوى يتطلب "تعديلًا متكررًا + تعاونًا متعدد الأشخاص" بشكل أفضل — من الأوراق البحثية إلى الوثائق، ومن الملاحظات إلى الخطط.
\end{itemize}

\end{tcbwarn}

\begin{tcbinfo}{بالنسبة لطلاب علوم الحاسوب، GitHub مهارة أساسية: إنها معرض أعمالك ومنصة تعاونك ومكتبة مصادر تعلمك، وأيضًا "سيرة ذاتية" مقنعة عند البحث عن عمل. سجّل حسابًا مبكرًا، وارفع واجباتك المقررية ومشاريعك الصغيرة، وتراكم ببطء، وبعد بضع سنوات ستملك سيرة تقنية تبهر المحاورين.}
جرّب: أكمل تجربتك الأولى في GitHub\textbf{سجّل حساب GitHub، وأكمل ملفك الشخصي (الصورة الرمزية، المقدمة، المجالات التقنية التي تهمك).}أنشئ مستودعًا عامًا باسم

\end{tcbinfo}

\textbf{hello-github}\par

، وحدد “Add a README file”.

\booksection{تاسعًا: الخلاصة — لماذا يجب على الجميع فهم GitHub}

في شريط بحث GitHub ابحث عن كلمة تهمك (مثل\textbf{python-tutorial})، و Star مشروعين يعجبانك.

(متقدم) ثبّت Git على الحاسوب، واستنسخ المستودع إلى المحلي، وعدّل README ثم ادفعه إلى السحابة، لإكمال سير عمل “محلي سحابي” كامل.

\begin{tcbpractice}{الفصل 24: تحليل متعمق لصلاحيات Root وقفل BL}
\begin{enumerate}
  \item هذا الفصل توسع متعمق لمحتوى الفصل 20، وهو مناسب للمستخدمين الذين يريدون فهم آليات الهاتف الأساسية فعلًا، أو حتى تجربة فلاش النظام. يُرجى ملاحظة أن هذه العمليات تنطوي على مخاطر معينة، فتصرف بحذر.
  \item أولًا: فهم متعمق لـ Bootloader\texttt{Bootloader هو أول برنامج يعمل عند إقلاع الهاتف، ومهمته تهيئة العتاد وتحميل نواة نظام التشغيل.\allowbreak  يمكن فهم Bootloader على أنه "حارس البوابة" للهاتف.\allowbreak }مبدأ عمل قفل BL:
  \item التحقق من التوقيع
  \item : يتحقق قفل BL من التوقيع الرقمي لصورة الإقلاع، فلا يقلع إلا الأنظمة الموقعة من الشركة المصنعة.\texttt{منع الرجوع للإصدار الأقدم}: يمنع المستخدم من تثبيت إصدارات أقدم لاستغلال الثغرات المعروفة.
  \item حماية البيانات
\end{enumerate}

\end{tcbpractice}
